\documentclass[12pt]{article}
\usepackage{amsmath, amssymb, amsfonts}
\usepackage{geometry}
\usepackage{hyperref}
\usepackage{graphicx}
\usepackage{bm}
\geometry{margin=1in}

\title{Wave Predictions Using the Kelvin Green's Function in the Limit $z\to 0^-$}
\author{G. D. Weymouth}
\date{\today}

\begin{document}

\maketitle

\begin{abstract}
Linear wave theory captures the essential physics of free-surface flows at a fraction of the computational cost of nonlinear and viscous methods, making it attractive for applications in design, real-time control, and surrogate modeling. However, linear wave predictions for ships with forward speed require evaluating the Kelvin Green's function at the free surface $z=0$, where the point-source kernel is ill-posed; the wave elevation grows without bound as $z\to 0^-$ even far downstream, causing both numerical and physical interpretation difficulties. In this paper we develop flat-ship theory for shallow-draft, high-speed planforms and show, via stationary-phase analysis, that its natural elliptic spanwise line integration acts as a wavenumber low-pass filter that exactly resolves this ill-posedness, yielding a kernel that is finite and differentiable at $z=0$. We then present a fast evaluator for both point and line kernels using contour deformation adapted to the non-analytic Kelvin phase, achieving $10^4$-$10^5$ speedup over direct quadrature while preserving far-wake asymptotics. An open-source Julia implementation is provided.
\end{abstract}

\section{Introduction}

Linear potential theory for steady forward-speed wave problems captures the essential physics of ship-wave interaction at a tiny fraction of the computational cost of nonlinear and viscous solvers. This computational efficiency makes it valuable for applications where many evaluations are required: parametric design optimization, real-time motion prediction for control systems, and as the physics-informed backbone for machine learning surrogate models.

The Kelvin Green's function for a steady translating disturbance exactly satisfies the linear free-surface boundary condition by construction. Peters (1949) established the integral representation for this function as a Rankine source-image pair and two additional terms: a smooth near-field integral and an oscillatory wavelike integral representing the radiated wave pattern (Noblesse 1981). For submerged disturbances, practical evaluation methods based on this approach have been available for decades including the fully analytic solution for a submerged spheroid by Farell (1978), and efficient Chebyshev surrogate representations for the smooth near-field term remain widely used for numerical approaches (Newman 1987).

However, the most direct output of linear wave theory---the free-surface elevation---is ill-posed in the point-source representation as $|z|\to 0^-$. In the classical Kelvin Green's function, the limit is non-uniform: the wave elevation spectrum $S_\zeta(k;z)$ peaks at $k^* \sim 1/|z|$ with amplitude $\sim 1/(|z|R)$, where $R$ is the planar distance from the source. As $|z|\to 0^-$ the peak migrates to arbitrarily high wavenumber while growing without bound, so no practical discretization can resolve the wavenumber content radiated over the extended downstream wake.

The problem is less severe for submeged bodies, but evaluation of the wavelike Green's function is still difficult for small submergence. Baar and Price (1988) developed a series representation, but the expansion is not uniformly convergent and does not have the correct asymptotic behaviour in the wake directly behind a disturbance. Iwashita and Ohkusu (1992) proposed a numerical steepest-descent approach, but unbounded wavenumbers and merging saddle points forced their method to be used well below $z=0$. Recent developments in numerical contour deformation (Huybrechs 2017; Gibbs \& Huybrechs 2024) provide a systematic framework for highly oscillatory integrals, but these methods are posed for analytic integrands, whereas the Kelvin wavelike integral has a non-analytic phase.

For surface-piercing bodies, the free-surface limit becomes unavoidable. Baar and Price showed that the waterline contour contributes directly to the potential through a line integral that arises when Green's second identity is applied in the fluid domain. This contour term can have a leading-order influence on the wave field, yet its evaluation requires the Green's function precisely at the free surface---the regime where existing representations fail. Alternative formulations have been proposed to circumvent the waterline contribution. The Neumann-Michell approach of Noblesse, Huang, and Yang (2013) replaces the explicit Green's function with an implicit iterative scheme that avoids the contour integral entirely. While this sidesteps the numerical difficulty, it sacrifices the directness and efficiency of an explicit kernel representation for the potential. The present work retains the classical formulation and addresses the modelling and numerical obstacle directly.

This paper addresses the $z\to 0^-$ limit with an emphasis on obtaining meaningful finite energy free-surface predictions. The first contribution is analytical: Section~\ref{sec:point_source_asymptotics} develops the limiting wave elevation spectrum $S_\zeta(k;z)\sim 1/(|z|R)$ stated above via stationary-phase analysis, showing that the point-source kernel is ill-posed; Section~\ref{sec:line_integrated_kernels} then developed the flat-ship model for shallow-draft $z\approx 0$ planforms, deriving the elliptic spanwise distribution of source strength that arises from the constant-downwash condition. The leading order flat-ship potential \textit{is} the waterline contribution, but analytic integration replaces the divergent spectrum with a finite $S_\zeta(k;0)\sim J_1(k)^2/k$ which decays as $k^{-3/2}$, where $J_1$ is the first-order Bessel function. The second contribution is numerical: Section~\ref{sec:fast_evaluator} develops a practical evaluator for both point and line kernels using contour deformation adapted to the non-analytic Kelvin phase, achieving $10^4$--$10^5$ speedup over direct quadrature while preserving the wave energy. The flat-ship theory, requiring only two (leading and trailing edge) kernel evaluations per field point, is then used to study the wave elevation and wave drag generated in the $z=0$ limit by a rectangular planform as a function of the length and beam based Froude-number. An open-source Julia implementation is provided for reproducibility.

\section{Mathematical Formulation}

\begin{figure}
    \centering
    \includegraphics[height=10cm]{example-image-a}
    \caption{Schematic of the flat ship model}
    \label{fig:placeholder}
\end{figure}

\subsection{Kelvin Green's Function}

Consider an inviscid, incompressible, irrotational flow in the half-space $Z<0$. The velocity potential $\Phi$ satisfies Laplace's equation
\[
\nabla^2\Phi = 0,\qquad Z<0,
\]
with the linearized free-surface boundary condition for steady forward speed $U$ in direction $X$:
\[
\ell\,\Phi_{XX} + \Phi_Z = 0,\qquad Z=0,
\]
where $\ell=U^2/g$ is the Kelvin length.

The Green's function at field point $\vec X$ due to a translating source at $\vec\alpha$ is
\[
G(\vec X,\vec\alpha) = -\frac 1{|\vec X-\vec\alpha|}+\frac 1{|\vec X-\vec\alpha'|}+\frac{N(\vec x)+W(\vec x)}\ell,
\]
where $\vec\alpha'$ is the image point reflected across $z=0$ and $\vec x = (x, y, z) = (\vec X-\vec\alpha')/\ell$ is the Kelvin-scaled vector from the reflected source.

The $N$ and $W$ functions are the near-field and wavelike terms, defined by the integral equations
\[
N = \frac 2\pi\int_{-1}^1 \Im\left(\text{expintx}\left((z\sqrt{1-T^2}+yT+i|x|)\sqrt{1-T^2}\right)\right) dT,
\]
and
\[
W = 4 H(-x)\int_{-\infty}^\infty \exp\bigl((1+t^2)z\bigr)\sin\bigl((x+|y|t)\sqrt{1+t^2}\bigr)\, dt,
\]
where $\text{expintx}(z)=e^z E_1(z)$ is the scaled exponential integral and $H$ is the Heaviside function ensuring the wave wake only propagates downstream. The near-field term is smooth and admits efficient Chebyshev polynomial representations. However, treatment of the wavelike term $W$ in the limit of $z\to 0^-$ is the central difficulty addressed here.

\subsection{Neumann-Kelvin Formulation}

In the Neumann-Kelvin formulation for steady forward speed, the potential is represented as a distribution of sources on the wetted surface of the body, plus a contribution from the waterline contour.

For a submerged body with wetted surface $S$ entirely below $Z=0$, the potential may be written as
\[
\Phi(\vec X) = \iint_S Q(\vec\alpha)\,G(\vec X,\vec\alpha)\,dS_{\alpha},
\]
where $Q$ is the source strength determined by enforcing the Neumann boundary condition $\partial_n\phi+Un_x=0$ on $S$ where $\hat n$ is the outward facing surface normal. The linear dependence on speed allows us to define the scaled strength $q=Q/U$ and scaled potential $\phi = \Phi/U\ell$. The scaled linear free surface elevation is then simply $\zeta = \partial_x\phi|_{Z=0}$.

For a surface-piercing body, let $C$ denote the waterplane footprint, with waterline contour $\partial C$. On the free surface $Z=0$, both $\phi$ and $G$ satisfy the linearized FSBC, which implies
\[
\Phi G_Z - G\Phi_Z = -\ell\,(\Phi G_{XX} - G\Phi_{XX}) = -\ell\,\frac{\partial}{\partial X}(\Phi G_X - G\Phi_X)
\]
Integrating over the exterior free surface and applying the divergence theorem reduces this to a contour integral over $\partial C$. Therefore, the potential for a surface-piercing body is a superposition of the submerged source distribution and the waterline contour contribution
\[
\phi(\vec X) = \ell\iint_S q(\vec\alpha)\,G(\vec X,\vec\alpha)\,ds_{\alpha}
\; + \; \ell\oint_{\partial C} q(\vec\alpha)\,G(\vec X,\vec\alpha)\,n_x\,dy_{\alpha}
\]
where $ds=dS/\ell^2$. Because the contour is on $Z=0$, the wave elevation $\zeta$ therefore requires evaluating the Green's function \textit{on} $z=0$. In the next section we show that the point-source kernel is ill-posed in this limit.

\section{Point-Source Asymptotics and Breakdown}\label{sec:point_source_asymptotics}

\begin{figure}
    \centering
    \includegraphics[height=6cm]{example-image-a}
    \includegraphics[height=6cm]{example-image-b}
    \caption{Stationary point locations as a function of $y/x$ for the point-source and line-integral after applying the Hankel function decomposition of equation~\ref{eq:Hankel}}
    \label{fig:placeholder}
\end{figure}

For a linearized velocity potential $\phi$ to yield physically meaningful predictions, both $\phi$ and its $x$-derivative $\partial_x\phi$ must be finite. This section analyzes the stationary-point structure of the oscillatory integral defining $W$ and shows that the dominant saddle migrates to $t_+ \to \infty$ as the wake centerline is approached, producing wave energy contributions that grow without bound. This motivates the need for regularization through spanwise integration, which is developed in Section \ref{sec:line_integrated_kernels}.

\subsection{A Unified Oscillatory-Integral Form}\label{ssec:unified}

In anticipation of the spanwise integration in Section \ref{sec:line_integrated_kernels}, we write the wavelike contribution in the general form
\begin{equation}
W_A(x,y,z) = 4H(-x)\int_{-\infty}^{\infty} A(t)\,\exp\left(z(1+t^2)\right)\,\sin\left(g(x,y,t)\right)\,dt,    
\end{equation}\label{eq:wavelike_general}
with phase $g(x,y,t) = (x + yt)\sqrt{1+t^2}$ and $A(t) = 1$ for the point-source case. The stationary points are given by $\partial_t g = 0$, which yields a quadratic equation in $t$:
\[
t_\pm = \frac{-x \pm \sqrt{x^2 - 8y^2}}{4y}.
\]
For $|y| > |x|/\sqrt{8}$ there are no real stationary points; for $|y| < |x|/\sqrt{8}$ there are two, corresponding to the transverse and diverging wave systems of the Kelvin wake. The critical observation is that in the near-centerline regime $|y| \ll |x|$, the larger root satisfies
\[
t_+ \approx -\frac{x}{2y}, \qquad |t_+| \to \infty \text{ as } y \to 0,
\]
so the dominant stationary point migrates to arbitrarily high values, corresponding to wavenumber $k_+ = t\sqrt{1 + t_+^2} \approx t_+^2$, as the centerline is approached.

\subsection{Point-Kernel Stationary-Phase Estimate}
\begin{figure}
    \centering
    \includegraphics[width=0.55\linewidth]{make_figures/wavecut.png}
    \includegraphics[width=0.44\linewidth]{make_figures/spectrum.png}
    \caption{Free surface wave cut and spectra downstream of a point-source with $z\to 0^-$ and line-integrate source on $z=0^-$. The wave cuts in the top and bottom plots are at $x=-8$ and $-40$, respectively. The line integrated results are for line half-beam $b=1$.}
    \label{fig:wavecut}
\end{figure}

A standard stationary-phase estimate at the large saddle $t_+$ uses $g_+'' \approx y_+$ and $y_+ \approx R/2t_+$ to give
\[
W_A(x,y,z) \sim A(t_+) \exp\left(-|z|\,t_+^2\right)\, \left(\frac{t_+}{R}\right)^{1/2}
\]
where $A \equiv 1$ for the point source. 

This result is problematic for linear wave theory for two reasons. First, it is clear that at $z=0$ there is no exponential damping, and the point-source produces unbounded wave energy: as $t_+ \to \infty$ the amplitude grows as $|W(x,y,0)| \sim (t_+/R)^{1/2}$. Second, even for finite submergence, while the exponential factor provides damping that suppresses the large-$t$ contributions, the peak wavenumber can easily become unresolvable for small $|z|$. In terms of $k$, the wavelike elevation spectrum $S^2_\zeta = |\partial_x W|^2$ for large $k$ at depth $z < 0$ is
%
\begin{equation}
  S^2_\zeta(k;\,z) \sim \frac{k}{R}\,\exp\left(-2|z|k\right),
  \label{eq:point_source_spectrum}
\end{equation}
%
which peaks at
%
\begin{equation}
  k^* = \frac{1}{2|z|}, \qquad S^2_\zeta(k^*;\,z) \;\sim\; \frac{1}{R\,|z|}.
  \label{eq:point_source_peak}
\end{equation}
%
As $|z| \to 0$ the peak migrates to arbitrarily high wavenumber while its amplitude grows as $|z|^{-1}$. For any fixed numerical resolution $k_{\max}$, the spectrum becomes unresolvable once $|z| \lesssim 1/(2k_{\max})$. The $R^{-1}$ decay along the wake means this is not a local problem: a near-surface source radiates an unresolvable wavenumber band over the entire downstream wake. Figure~\ref{fig:wavecut} shows wave cuts and their spectra downstream of a point source for a sequence of decreasing $|z|$ values, illustrating the peak migration and amplitude growth.

The natural regularization in the waterline formulation is spanwise line integration, developed in the next section.

\section{Line-Integrated Potential and Flat-Ship Theory}\label{sec:line_integrated_kernels}

This section develops the line-integrated Kelvin kernel that arises naturally in waterline-driven configurations. We first establish the flat-ship model as the physical context, derive the elliptic spanwise distribution, and then show that the resulting line-integrated kernel is finite and differentiable at $z=0$ --- in sharp contrast to the point-source behavior analyzed in Section~\ref{sec:point_source_asymptotics}.

\subsection{The Flat-Ship Model}\label{ssec:flat_ship_model}

Consider a horizontal surface-piercing planform, e.g., a planing hull, at small angle of attack $\alpha$. The flat-ship idealization imposes uniform downwash approaching the free surface
\[
\Phi_Z = \alpha U,\qquad (X,X)\in S,\ Z\to 0^-.
\]
For a rectangular planform of half-width $b\ell$ and length $L\ell$, the potential reduces to contributions from the leading and trailing edges:
\[
\phi(\vec X) = \ell\int_{-b}^{b} q(y')\left[G(x,y-y',z)-G(x-L,y-y',z)\right]dy'.
\]
The remaining unknown is the spanwise distribution $q(y')$. For uniform downwash forcing, the dominant nearfield operator is logarithmic in span:
\[
\int_{-b}^{b} q(y')\log|y-y'|\,dy' = \mathrm{const}.
\]
This is the classical constant-downwash problem from finite-wing theory and the solution on $[-b,b]$ is the elliptic distribution
\[
q(y') = q_0\sqrt{1-(y'/b)^2}.
\]

\subsection{The Line-Integrated Wavelike Kernel}

We define the integrated wavelike contribution along a spanwise edge as
%
\begin{align}
  W_b(x,y,z) &= \int_{-b}^{b} \sqrt{1-(y'/b)^2}\,W(x,y-y',z)\,dy' \\
  &=  4H(-x)\int_{-\infty}^{\infty} \exp\left(z(1+t^2)\right) \int_{-b}^{b} \sqrt{1-(y'/b)^2} \,\sin\left(g(x,y,t)\right),dy'\,dt,
  \label{eq:Wb_def}
\end{align}
%
where we have swapped the order of integration. This is the Fourier transform of the elliptic distribution which has a known Bessel function representation, giving
\begin{equation}
  W_b(x,y,z) = 4\pi H(-x)\int_{-\infty}^{\infty} A_b(t)\,\exp\left(z(1+t^2)\right)\,\sin\left(g(x,y,t)\right)\,dt,
  \label{eq:Wb_integral}
\end{equation}
%
with
%
\begin{equation}
  A_b(t) = \pi\,\frac{J_1\bigl(b\,k(t)\bigr)}{k(t)},\qquad k(t)=t\sqrt{1+t^2}.
  \label{eq:Ab}
\end{equation}
%
The factor $J_1(bk)/k$ is finite at $k=0$ (with limit $b/2$) and decays as $k \to \infty$. This is the mechanism by which spanwise smoothing regularizes the kernel: the elliptic distribution acts as a low-pass filter in wavenumber, suppressing the large-$k$ contributions that drove the point-source divergence in Section~\ref{ssec:unified}.

The wave elevation spectrum at $z=0$ is
%
\begin{equation}
  S_{b,\zeta}(k;\,0) \sim \frac{J_1^2\left(b\,k(t)\right)}{k}
  \label{eq:Sb_spectrum}
\end{equation}
%
which is $O(b^2)$ at $k=0$ and decays as $k^{-2}$ for large $k$ --- integrable, with a finite peak set by the beam $b$. This completely eliminates the divergent high wavenumber energy found in the point-source as shown in Figure~\ref{fig:wavecut} for $b=1$.

\subsection{Wave Resistance}
\label{sec:wave_resistance}

The flat-ship wave resistance $R_w$ follows from the Havelock formula. Working directly in the variable $t = \sin\theta/\cos^2\theta$ with Jacobian $d\theta/\cos^3\theta = (1+t^2)/(1+2t^2)\,dt$, the Havelock formula becomes
%
\begin{equation}
  R_w = \frac{\rho}{2\pi\ell^2} \int_{-\infty}^{\infty} |K(t)|^2\, \frac{1+t^2}{1+2t^2}\, dt,
  \label{eq:havelock}
\end{equation}
%
where $K$ is the Kochin function, the source-strength-weighted integral of the wavelike kernel over the planform. Defining the wave drag coefficient $C_W = R_w/(\rho U^2 b^2\ell^2)$ and substituting the elliptic spanwise distribution derived above gives
%
\begin{equation}
  C_W = \pi q_0^2 \int_{-\infty}^{\infty}
  A_w(t)\bigl(1 - \cos\bigl(L\sqrt{1+t^2}\bigr)\bigr)\, dt\quad \text{where}\quad A_w(t)=\frac{J_1(bk(t))^2}{b^2t^2(1+2t^2)}.
  \label{eq:Cw}
\end{equation}
%
with $q_0 \propto \alpha$. The integrand is $O(1)$ at $k=0$ and decays as $k^{-3}$ for large $k$, so~\eqref{eq:Cw} is well-posed and rapidly convergent. The Bessel factor $(J_1(bk)/bk)^2$ controls angular energy spreading with beam $b$ while the interference factor $1-\cos(L\sqrt{1+t^2})$ oscillates with $L$, producing the classical constructive and destructive interference between bow and stern wave systems.

\section{Unified contour deformation framework} \label{sec:fast_evaluator}

The general wavelike integral of equation~\ref{eq:wavelike_general}, which covers the point-source, flat-ship, and Havelock cases, is evaluated using a unified contour-deformation framework following the Path-Finder approach of Gibbs and Huybrechs (2024). The key idea is to partition the integral into non-oscillatory intervals around stationary points, where standard quadrature is efficient, and deform the remaining semi-infinite tails into the complex plane, where the integrand decays exponentially. This utilizes the saddle point analysis of the previous sections to develop a uniformly valid, accurate, and efficient evaluation method for the wavelike terms.

\subsection{Defining the saddle regions and quadrature methods}

\begin{figure}
    \centering
    \includegraphics[width=0.495\linewidth]{make_figures/dphase_z.png}
    \includegraphics[width=0.495\linewidth]{make_figures/dphase_GLn.png}
    \caption{Error convergence of the partitioned quadrature method applied to the same wavecuts in figure \ref{fig:wavecut}.}
    \label{fig:delta_phase}
\end{figure}

Unlike the analytic phase functions treated in Gibbs and Huybrechs (2024) the wavelike phase $g(x,y,t) = (x+yt)\sqrt{1+t^2}$ introduces branch points at $t = \pm i$. To avoid these points, we limit the non-oscillatory intervals to the real axis such that each endpoint is separated from the nearest stationary point by a prescribed phase increment $\Delta g$. The non-analytic phase also means that these partition boundaries must be located by numerical 1D root-finding. 

The kernel is smooth over the non-oscillatory interval by construction, meaning a Gauss-Legendre quadrature with a fixed number of points related to $\Delta g$ can be applied. For this manuscript, we instead use an adaptive Gauss-Konrad quadrature to ensure a prescribed accuracy. The semi-infinite tails are evaluated using numerical steepest descent on a complex path emanating from the interval endpoints using Gauss-Laguerre quadrature, which is appropriate for exponentially decaying kernels. The contour points are located by Newton's method, selecting the $\sqrt{1+t^2}$ brand such that the phase remains continuous along each contour. 

Figure~\ref{fig:delta_phase} shows the maximum error of the point-source wavelike integral using the partioning method described above for the same wave cuts as in Figure~\ref{fig:wavecut}. The error envelop is insensitive to $\vec x$, handling the wave cusp merge as $y\to \pm x/\sqrt{8}$ and the energy growth as $z\to 0^-$ gracefully. The phase spacing $\Delta g$ is the primary numerical parameter of the approach, encoding a tradeoff between the number of oscillations integrated along the real line (and therefore the difficulty of that integration) and the exponential decay along the contours (which determines the accuracy of that quadrature).  The secondary numerical parameter is the number of Gauss-Laguerre points $N$ used on each contour. Since each point must be found using an iterative Newton search, fewer points are preferred, and we use $\Delta g \approx 6$ and $N=4$ producing maximum errors around $10^{-6}$.

\subsection{Hankel function decomposition}

\begin{table}[]
    \centering
    \begin{tabular}{c|c}
         &  \\
         & 
    \end{tabular}
    \caption{Caption}
    \label{tab:wavelike}
\end{table}

The same contour-deformation strategy applies to the line-integrated kernel, extended to handle the additional oscillatory structure introduced by the Bessel prefactor $J_1(bk(t))/k(t)$. The Bessel function is decomposed using Hankel functions away from $t = 0$
\begin{equation}
J_1(z) = \frac{1}{2}\left(H_1^{(1)}(z)\,e^{iz} + H_1^{(2)}(z)\,e^{-iz}\right)    
\end{equation}\label{eq:Hankel}
The scaled Hankel functions $H_1^{(1,2)}$ are slowly varying and the exponential factors can be absorbed into the complex phase. The line-integrate kernel is therefore expressed as a sum of two oscillatory integrals with shifted phases $g(x,y \pm b)$. The same partitioning strategy is applied to each phase function, and their union defines full non-oscillatory range. This region is integrated once with the Bessel prefactor weighting, and the semi-infinite tails are computed twice, once for each Hankel function using the appropriate phase function. 

Overall, the line-integral evaluation is around twice as expensive as the point-source evaluation, but Table\ref{tab:wavelike} shows provides $10^4$--$10^5$ speedup over direct quadrature on $z=0$ with relative errors less than $10^{-6}$. Because the flat-ship limit reduces the planform influence to just two spanwise line-integrals, two calls this integral function are sufficient to evaluate the potential at any field point. 

Note that we can use the same approach for the wave resistance integral~\eqref{eq:Cw}. Applying the Hankel function decomposition to $J_1^2$ results in three terms, two that oscillate as $e^{\pm2bk}$ and one slowly varying part, meaning we would need to handle three phase functions: $g(x,y,t),\,g(x,y\pm2b,t)$. The fast $t^{-6}$ decay of $A_w$ means direct quadrature is also practical. 

\section{Wavefield and wave drag predictions}

\begin{figure}
    \centering
    \includegraphics[width=0.495\linewidth]{make_figures/wave_field_b1.png}
    \includegraphics[width=0.495\linewidth]{make_figures/wave_field_b2.png}
    \caption{Wave elevation $\zeta$ scaled by the spanwise integrated strength $\tfrac \pi 2 q_0b$ for $b=1,\, 2$ on the left and right, respectively.}
    \label{fig:placeholder}
\end{figure}

\section*{References (incomplete)}

\begin{itemize}
    \item Baar, J. J. M. and Price, W. G. (1988). \textbf{TODO} (full bibliographic entry).
    \item Noblesse, F. (1981). \textbf{TODO} (full bibliographic entry).
    \item Newman, J. N. (1987). \textbf{TODO} (full bibliographic entry).
    \item Peters, A. S. (1949). \textbf{TODO} (full bibliographic entry).
    \item Iwashita, H. and Ohkusu, M. (1992). \textbf{TODO} (full bibliographic entry).
    \item Noblesse, F., Huang, F., and Yang, C. (2013). \textbf{TODO} (full bibliographic entry; Neumann-Michell reformulation).
    \item Huybrechs, D. (2017). \textbf{TODO} (full bibliographic entry).
    \item Gibbs, A. and Huybrechs, D. (2024). \textbf{TODO} (full bibliographic entry).
\end{itemize}

\end{document}